\section{Interior filling and interpolation}\label{sec:interpolation}

The reshaping construction preserves the number of exponents, but does not
identify the individual exponents that survive.  Gradedness supplies the missing
information: a family with full asymptotic density contains every lattice point
in a fixed compact subset of the interior, in all sufficiently large degrees.
We first prove this statement, then use it to find a complete interpolation
space inside the reshaped supports.  These two arguments follow
\cite[Lemmas~5.1 and~5.2]{OpenAI2026Higher}.  Interior filling for semigroups
is studied more generally in \cite[Theorem~1.6]{KavehKhovanskii2012}; the
full-density hypothesis here permits a direct proof by counting decompositions.

\begin{lemma}[Interior filling]\label{lem:interior-filling}
Let $P\subset\R^n$, $n\ge1$, be a compact convex set with nonempty interior.
Suppose finite sets $F_N\subset NP\cap\Z^n$, indexed by integers $N\ge1$,
satisfy
\[
 F_N+F_{N'}\subset F_{N+N'}
 \quad\text{and}\quad
 \#F_N=\vol(P)N^n+o(N^n).
\]
Then for every compact set $K\subset\operatorname{int}(P)$ there is an integer
$N_K$ such that
\[
 NK\cap\Z^n\subset F_N\qquad(N\ge N_K).
\]
\end{lemma}

\begin{proof}
The boundary of a convex body has measure zero, so lattice-point counting by
Riemann sums gives
\[
 \#(NP\cap\Z^n)=\vol(P)N^n+o(N^n).
\]
Consequently the number of missing lattice points
\[
 b_N:=\#\bigl((NP\cap\Z^n)\setminus F_N\bigr)
\]
is $o(N^n)$.  We show that a point of $NK$ has too many decompositions into two
points in nearly equal degrees for all of them to involve a missing point.

Assume $K$ is nonempty, and choose $\rho>0$ so that the closed ball of radius
$\rho$ about every point of $K$ lies in $P$.  Write
$N=N_1+N_2$ with $N_1=\lfloor N/2\rfloor$ and
$N_2=\lceil N/2\rceil$.  For $y\in NK\cap\Z^n$, consider all lattice points
\[
 y_1\in\Z^n,
 \qquad
 \left\|y_1-\frac{N_1}{N}y\right\|<\rho N_1,
 \qquad y_2:=y-y_1.
\]
Both $y_1/N_1$ and $y_2/N_2$ lie in the radius-$\rho$ ball about $y/N$.
Thus $y_i\in N_iP\cap\Z^n$.  The ball from which $y_1$ is chosen contains
at least $cN^n$ lattice points for a constant $c>0$ independent of $y$ and
all sufficiently large $N$: an inscribed cube gives this bound uniformly in
the center.  Among these choices, at most $b_{N_1}$ have $y_1\notin F_{N_1}$,
and at most $b_{N_2}$ have $y_2\notin F_{N_2}$, since $y_1\mapsto y-y_1$ is
injective.  As $b_{N_1}+b_{N_2}=o(N^n)$, at least one choice has
$y_i\in F_{N_i}$.  Gradedness then gives $y=y_1+y_2\in F_N$.
All bounds are uniform in $y$, which proves the assertion.
\end{proof}

The interpolation space we need is a simplex stretched in the last coordinate
by the number of points.  Its fibers over the other exponents have exactly the
lengths needed for one-variable interpolation with multiplicities.

\begin{lemma}[Jets from an elongated simplex]\label{lem:elongated-jets}
Let $n,r\ge1$ and $m\ge0$ be integers, and let $k$ be a field containing
distinct nonzero elements $c_1,\ldots,c_r$.  Set
\begin{equation}\label{eq:interpolation-support}
 E_m=\left\{(\beta,b)\in\Z_{\ge0}^{n-1}\times\Z_{\ge0}:
          |\beta|+\frac br<m+1\right\},
 \qquad |\beta|=\sum_{j=1}^{n-1}\beta_j.
\end{equation}
Then $M(E_m)$ separates jets through order $m$ at the $r$ points
\[
 q_i=(1,\ldots,1,c_i)\in\T^n(k).
\]
For $n=1$, the tuple $\beta$ is empty and $|\beta|=0$.
\end{lemma}

\begin{proof}
Write the coordinates as $(x_1,\ldots,x_{n-1},Y)$ and put
$x'-1=(x_1-1,\ldots,x_{n-1}-1)$.  An arbitrary jet through order $m$ at $q_i$
has a unique expression
\[
 \sum_{|\alpha|\le m}(x'-1)^\alpha h_{i,\alpha}(Y-c_i),
 \qquad \deg h_{i,\alpha}\le m-|\alpha|.
\]
For each $\alpha$, the Chinese remainder theorem gives a polynomial
$h_\alpha(Y)$ satisfying
\[
 h_\alpha(Y)\equiv h_{i,\alpha}(Y-c_i)
   \pmod{(Y-c_i)^{m-|\alpha|+1}}
 \quad(1\le i\le r),
 \qquad
 \deg h_\alpha<r(m-|\alpha|+1).
\]
Indeed, the ideals in this display are pairwise comaximal, and their product
is generated by a monic polynomial of degree $r(m-|\alpha|+1)$.
The polynomial
\[
 \sum_{|\alpha|\le m}(x'-1)^\alpha h_\alpha(Y)
\]
realizes the prescribed jets.  Each monomial $x'^\beta Y^b$ in its expansion
satisfies $\beta_j\le\alpha_j$ for every $j$ and
$b<r(m-|\alpha|+1)$ for some $\alpha$.  Hence
$|\beta|+b/r<m+1$, so the polynomial belongs to $M(E_m)$.
This uses powers of local parameters and the Chinese remainder theorem,
and therefore works in every characteristic.
\end{proof}

We now choose the number of points large enough that successive applications
of Proposition~\ref{prop:reshape} produce the required elongated simplex.

\begin{proposition}[Jets at the geometric generic tuple]\label{prop:generic-jets}
Let $k$ be an algebraically closed field of positive characteristic, let $X/k$
be a smooth integral projective variety of dimension $n\ge2$, and let $L$ be
ample.  Put $V=L^n$.  There is an integer $r_0=r_0(X,L)$ such that, for every
integer $r\ge r_0$ and every real number
\[
 0<\theta<\left(\frac Vr\right)^{1/n},
\]
there is an integer $N_0=N_0(r,\theta)$ such that, for every $N\ge N_0$,
the sections of $NL$ separate jets through order $\lfloor N\theta\rfloor$
at some ordered $r$-tuple of distinct $k$-points, which may depend on $N$.
For each such $N$, the same jet separation holds after base change to $K_r$
at the geometric generic tuple $p_1,\ldots,p_r$.
\end{proposition}

\begin{proof}
Take the graded family $S_N$ and simplex
$P_0=\Delta(a_1,\ldots,a_n)$ from Proposition~\ref{prop:initial-simplex}.
Thus $\prod_j a_j=V$, the support $S_N$ lies in $NP_0$, and
\[
 \#S_N=\frac{V}{n!}N^n+o(N^n).
\]
Let $\ell=\min_j a_j$, and choose $r_0$ so large that
\[
 w:=\left(\frac Vr\right)^{1/n}<\frac{\ell}{100}
 \qquad(r\ge r_0).
\]
Fix one such $r$.  Reshape the first two intercepts to $(w,a_1a_2/w)$,
then reshape the second and third intercepts, and continue through the
$(n-1)$st and $n$th intercepts.  At every step, both intercepts $A,B$
being reshaped are at least $\ell$: untouched intercepts have this property,
and the carried intercept $AB/w$ is larger than $A$ because $B>w$.
Consequently
\[
 w\max(3/A,2/B)\le\frac{3w}{\ell}<\frac1{16},
\]
as required by Proposition~\ref{prop:reshape}.  The product of all intercepts
remains $V$.  The resulting graded family $F_N$ therefore satisfies
\begin{equation}\label{eq:elongated-full-density}
 \begin{aligned}
 &F_N\subset NP_*\cap\Z^n,
 \qquad P_*:=\Delta(w,\ldots,w,rw),\\
 &\#F_N=\frac{V}{n!}N^n+o(N^n)
        =\vol(P_*)N^n+o(N^n).
 \end{aligned}
\end{equation}
Here the last intercept is $V/w^{n-1}=rw$.
The transformations preserve backward jet transfer, so jet separation by
$M(F_N)$ will imply jet separation by $H^0(X,NL)$.

Fix $0<\theta<w$.  To apply Lemma~\ref{lem:interior-filling}, we contract
$P_*$ and then translate it away from the coordinate faces.
Choose $\lambda$ with $\theta/w<\lambda<1$, and choose a vector
$v\in\R_{>0}^n$ small enough that
\[
 v+\lambda P_*\subset\operatorname{int}(P_*).
\]
For example, it suffices that
$\sum_{j<n}v_j/w+v_n/(rw)<1-\lambda$.
Choose vectors $v_N\in N^{-1}\Z^n$ tending to $v$.
For all sufficiently large $N$, the sets $v_N+\lambda P_*$ lie in a common
compact subset $K$ of $\operatorname{int}(P_*)$.  Applying
Lemma~\ref{lem:interior-filling} to \eqref{eq:elongated-full-density} gives
\begin{equation}\label{eq:translated-inner-support}
 Nv_N+\bigl(N\lambda P_*\cap\Z^n\bigr)\subset F_N
\end{equation}
for all sufficiently large $N$.

Put $m=\lfloor N\theta\rfloor$.  Since $\theta<\lambda w$, we have
$m+1\le N\lambda w$ for all sufficiently large $N$.
The defining inequality \eqref{eq:interpolation-support} therefore gives
\[
 E_m\subset N\lambda P_*\cap\Z^n.
\]
Choose distinct $c_1,\ldots,c_r\in k^\times$ and apply
Lemma~\ref{lem:elongated-jets}.  Multiplication by the Laurent monomial with
exponent $Nv_N\in\Z^n$ is multiplication by a unit on the torus, hence
induces an automorphism on every jet space.  Thus
\eqref{eq:translated-inner-support} implies that $M(F_N)$ separates jets
through order $m$ at these $r$ points.  Backward jet transfer through all
reshapings and Proposition~\ref{prop:initial-simplex} gives the asserted
jet separation at an ordered tuple on $X$.

For each fixed $r$, $\theta$, and sufficiently large $N$, the locus in $U_r$
where this evaluation map is surjective is a nonempty open set, by the
openness of jet separation established in Section~\ref{sec:transfer}.
The variety $U_r$ is integral, so its generic point belongs to that open set.
Base change of the jet evaluation map to $K_r$ preserves surjectivity and
identifies its targets with the jets at $p_1,\ldots,p_r$.  The same
geometric generic tuple therefore satisfies every required jet-separation
statement.
\end{proof}
